\chapter{Эх кодын жагсаалт}

Энэхүү хавсралтад системийн хөгжүүлэлтэд ашигласан голлох эх кодуудыг хавсаргав. 

\section{Үндсэн API сервер (api\_server.py)}
Системийн backend хэсэг буюу FastAPI дээр бичигдсэн үндсэн сервер.
\begin{verbatim}
from fastapi import FastAPI, BackgroundTasks, HTTPException, Depends
from pydantic import BaseModel
import uvicorn
import os
import json
import subprocess
from datetime import datetime
from sqlalchemy.orm import Session
from sqlalchemy import func

from db.database import DatabaseManager, FacebookPost
from ml.analyzer import DataAnalyzer
from ml.gemini_integration import GeminiAnalyzer
from ml.sheets_integration import SheetsExporter

app = FastAPI(title="Social Media Analysis API", version="1.0.0")
db_manager = DatabaseManager()
analyzer = DataAnalyzer()
gemini = GeminiAnalyzer()

@app.get("/health")
def health_check():
    return {
        "status": "ok",
        "database": db_manager.engine is not None,
        "analyzer": analyzer.is_ready,
        "gemini": gemini.is_ready
    }

@app.get("/api/v1/stats")
def get_stats():
    session = db_manager.get_session()
    try:
        total_posts = session.query(FacebookPost).count()
        total_comments = session.query(func.sum(FacebookPost.comments)).scalar() or 0
        avg_engagement = session.query(func.avg(
            FacebookPost.likes + FacebookPost.comments + FacebookPost.shares
        )).scalar() or 0
        
        return {
            "status": "success",
            "data": {
                "total_posts": total_posts,
                "total_comments": int(total_comments),
                "avg_engagement": round(float(avg_engagement), 2)
            }
        }
    except Exception as e:
        raise HTTPException(status_code=500, detail=str(e))
    finally:
        session.close()

@app.post("/api/v1/scraper/run")
def run_scraper(background_tasks: BackgroundTasks):
    try:
        subprocess.Popen(["python", "run_scraper.py", "--headless", "--max-posts=10"])
        return {"status": "success", "message": "Scraper started in background"}
    except Exception as e:
        raise HTTPException(status_code=500, detail=str(e))
\end{verbatim}

\section{React Dashboard (App.tsx)}
Системийн frontend интерфэйсийг удирддаг үндсэн React компонент.
\begin{verbatim}
import React, { useState, useEffect } from 'react';
import { BrowserRouter, Routes, Route } from 'react-router-dom';
import Sidebar from './components/Sidebar';
import Dashboard from './components/Dashboard';
import DataSources from './components/DataSources';
import AIAnalysis from './components/AIAnalysis';
import Settings from './components/Settings';

function App() {
  const [stats, setStats] = useState(null);

  useEffect(() => {
    fetch('/api/v1/stats')
      .then(res => res.json())
      .then(data => setStats(data.data))
      .catch(err => console.error("Failed to load stats", err));
  }, []);

  return (
    <BrowserRouter>
      <div className="flex h-screen bg-gray-50">
        <Sidebar />
        <main className="flex-1 overflow-y-auto p-8">
          <Routes>
            <Route path="/" element={<Dashboard stats={stats} />} />
            <Route path="/sources" element={<DataSources />} />
            <Route path="/analysis" element={<AIAnalysis />} />
            <Route path="/settings" element={<Settings />} />
          </Routes>
        </main>
      </div>
    </BrowserRouter>
  );
}
export default App;
\end{verbatim}

\section{ML Analyzer (analyzer.py)}
Өгөгдөлд NLP шинжилгээ хийдэг гол класс.
\begin{verbatim}
import pandas as pd
from textblob import TextBlob
from sklearn.feature_extraction.text import CountVectorizer
from sklearn.decomposition import LatentDirichletAllocation
import re

class DataAnalyzer:
    def __init__(self):
        self.is_ready = True
        
    def analyze_sentiment(self, text):
        if not text:
            return 0.0
        blob = TextBlob(str(text))
        return blob.sentiment.polarity

    def extract_topics(self, documents, n_topics=3):
        if len(documents) < 5:
            return ["Insufficient data for topic modeling"]
            
        vectorizer = CountVectorizer(stop_words='english', max_features=1000)
        doc_term_matrix = vectorizer.fit_transform(documents)
        
        lda = LatentDirichletAllocation(n_components=n_topics, random_state=42)
        lda.fit(doc_term_matrix)
        
        feature_names = vectorizer.get_feature_names_out()
        topics = []
        for topic_idx, topic in enumerate(lda.components_):
            top_words = [feature_names[i] for i in topic.argsort()[:-6:-1]]
            topics.append(f"Topic {topic_idx+1}: {', '.join(top_words)}")
            
        return topics
\end{verbatim}
